\documentclass[12pt]{article}

\usepackage{sbc-template}
\usepackage{graphicx,url}
\usepackage[utf8]{inputenc}
\usepackage[T1]{fontenc}
\usepackage{xcolor}
\usepackage{booktabs}
\usepackage{tabularx}
\usepackage{array}
\IfFileExists{portuges.ldf}{\usepackage[brazil]{babel}}{%
  \renewcommand{\figurename}{Figura}\renewcommand{\tablename}{Tabela}%
  \renewcommand{\refname}{Referências}}

\sloppy

% Marcador de conteúdo que depende da execução do notebook
\newcommand{\pend}[1]{\textcolor{red}{\textbf{[#1]}}}
% Figura gerada pelo notebook (mostra aviso enquanto o arquivo não existir)
\newcommand{\figura}[4]{%
\begin{figure}[ht]\centering
\IfFileExists{figuras/#1}{\includegraphics[width=#2\textwidth]{figuras/#1}}%
{\fbox{\parbox[c][3cm][c]{#2\textwidth}{\centering\pend{figura \texttt{\detokenize{#1}} gerada pelo notebook}}}}
\caption{#3}\label{#4}
\end{figure}}

\title{Agente Identificador de Fraudes Bancárias: Detecção de Fraudes em
Transações com Cartão de Crédito por Aprendizado de Máquina Supervisionado}

\author{Cheuk Ki Yu\inst{1}, Gabriela Nellessen de Sousa\inst{1}, Milton Almeida Leoncio\inst{1}}

\address{Faculdade de Computação e Informática -- Universidade Presbiteriana Mackenzie\\
  São Paulo -- SP -- Brasil
  \email{\{10419664, 10441930, 10416764\}@mackenzista.com.br}
}

\begin{document}

\maketitle

\begin{abstract}
Card fraud imposes large losses on consumers and financial institutions, and
its detection is hindered by extreme class imbalance. This partial report
formalizes fraud detection as a supervised binary classification task on the
public ULB Credit Card Fraud dataset (284,807 transactions, 0.172\% frauds).
We present an exploratory data analysis, a leakage-free preparation pipeline
(deduplication, cyclic time encoding, log-scaled amount, stratified and
temporal splits, SMOTE inside the training folds) and baseline results for
Logistic Regression and Random Forest, evaluated with PR-AUC, recall and MCC
instead of accuracy. The best baseline, a class-weighted Random Forest, reached PR-AUC of 0.80 and MCC of 0.84, while accuracy proved uninformative (99.8\% for a trivial classifier). The next
stage will compare tuned Random Forest, XGBoost and SVM, add explainability
and discuss ethical risks such as false positives.
\end{abstract}

\begin{resumo}
Fraudes com cartão geram grandes perdas a consumidores e instituições
financeiras, e sua detecção é dificultada pelo extremo desbalanceamento entre
classes. Este artigo parcial formaliza a detecção de fraudes como classificação
binária supervisionada sobre o dataset público Credit Card Fraud (ULB), com
284.807 transações e 0,172\% de fraudes. São apresentadas a análise
exploratória, um processo de preparação sem vazamento de dados (remoção de
duplicatas, codificação cíclica do tempo, transformação logarítmica do valor,
divisões estratificada e temporal e SMOTE apenas no treino) e resultados de
referência com Regressão Logística e Random Forest, avaliados por PR-AUC,
Recall e MCC em vez de acurácia. O melhor modelo de referência, uma Random Forest com pesos balanceados, obteve
PR-AUC de 0,80 e MCC de 0,84, enquanto a acurácia se mostrou pouco informativa (99,8\% para um classificador trivial). A próxima etapa comparará Random Forest ajustada, XGBoost e
SVM, com explicabilidade e análise de riscos éticos.
\end{resumo}

% =====================================================================
\section{Introdução}

\subsection{Contextualização}

A digitalização dos meios de pagamento ampliou o volume de transações
eletrônicas e, com ele, as oportunidades de fraude. No Brasil, a Federação
Brasileira de Bancos (Febraban) estimou em R\$~10,1 bilhões as perdas com
golpes e fraudes em canais eletrônicos e cartões em 2024, alta de 17\% em
relação a 2023 \cite{poder360:2025}. Para as instituições financeiras,
identificar rapidamente uma transação fraudulenta entre milhões de transações
legítimas é um problema de decisão sob incerteza, no qual regras fixas escritas
por especialistas tendem a ficar desatualizadas diante de novas estratégias dos
fraudadores \cite{dalpozzolo:2018}.

O aprendizado de máquina oferece uma alternativa: modelos que aprendem, a partir
de transações históricas rotuladas, padrões que distinguem operações legítimas
de fraudulentas. O problema, porém, tem uma dificuldade central: as fraudes são
raras. No conjunto de dados usado neste trabalho, elas correspondem a cerca de
0,17\% das transações \cite{kaggle:ulb}, o que torna métricas usuais, como a
acurácia, enganosas e exige técnicas específicas de tratamento do
desbalanceamento \cite{chawla:2002,saito:2015}.

\subsection{Justificativa}

Um sistema de detecção precisa equilibrar dois tipos de erro com custos
distintos. O \emph{falso negativo} (fraude não detectada) gera prejuízo
financeiro direto. O \emph{falso positivo} (transação legítima bloqueada) gera
atrito com o cliente, custo operacional de análise manual e perda de confiança.
Comparar de forma controlada diferentes algoritmos e estratégias de
balanceamento, com métricas adequadas a classes raras, fornece evidência
quantitativa para escolher modelos que reduzam ambos os erros. O tema se alinha
aos Objetivos de Desenvolvimento Sustentável (ODS) 8, por contribuir para um
sistema financeiro mais seguro e com menos perdas econômicas; 9, pela aplicação
de ciência de dados e inovação tecnológica à infraestrutura financeira; e 16,
pelo combate a crimes financeiros.

\subsection{Motivação}

A motivação do grupo está em aplicar os conteúdos da disciplina a um problema
real, de alto impacto e com dificuldades técnicas não triviais (desbalanceamento
extremo, atributos anonimizados e custo assimétrico dos erros). O projeto também
será aprofundado no Trabalho de Conclusão de Curso dos integrantes, o que
reforça a importância de uma base metodológica sólida e reprodutível desde esta
etapa.

\subsection{Objetivo}

Desenvolver e avaliar modelos de aprendizado de máquina supervisionado para
classificar transações com cartão de crédito como legítimas ou fraudulentas,
comparando algoritmos (Regressão Logística como referência, Random Forest,
XGBoost e SVM) e estratégias de tratamento do desbalanceamento (pesos de classe
e SMOTE), com ênfase na maximização da detecção de fraudes (Recall) e no
controle dos falsos positivos (Precisão).

\subsection{Contribuições esperadas}

Espera-se contribuir com: (i) a formalização da detecção de fraudes como tarefa
de classificação binária supervisionada; (ii) uma análise exploratória
documentada do dataset, identificando suas características e limitações; (iii)
um protocolo experimental sem vazamento de dados, com validação estratificada e
temporal; (iv) a comparação quantitativa entre modelos e estratégias de
balanceamento por métricas adequadas a classes raras; (v) a análise dos
atributos mais relevantes às decisões dos modelos; (vi) um protótipo capaz de
classificar novas transações; e (vii) um repositório público que permita
reproduzir todos os experimentos.

\subsection{Evolução em relação à proposta}

Em relação à proposta inicial, o escopo foi refinado em quatro pontos. Primeiro,
o dataset foi definido (Seção~\ref{sec:dataset}). Segundo, a acurácia deixou de
ser métrica de decisão, pois um classificador que rotula todas as transações
como legítimas já atinge cerca de 99,8\%; as métricas principais passaram a ser
PR-AUC, Recall, Precisão, F1 e MCC \cite{saito:2015,chicco:2020}. Terceiro, foi
acrescentada uma validação com divisão temporal, mais próxima do uso real. Por
fim, como o dataset cobre apenas dois dias, a adaptação a novas estratégias de
fraude ao longo do tempo (\emph{concept drift}) será discutida como limitação,
e não avaliada experimentalmente.

% =====================================================================
\section{Fundamentação Teórica}

\subsection{Agentes, aprendizado supervisionado e classificação}

Em Inteligência Artificial, um agente é uma entidade que percebe o ambiente e
age sobre ele \cite{russell:2021}. No sistema proposto, o agente percebe cada
nova transação, estima sua probabilidade de fraude por meio de um modelo
aprendido e age emitindo ou não um alerta para análise. O modelo é obtido por
\emph{aprendizado supervisionado}: dado um conjunto de exemplos
$\{(\mathbf{x}_i, y_i)\}$, em que $\mathbf{x}_i$ é o vetor de atributos da
transação e $y_i \in \{0, 1\}$ indica se ela é fraudulenta, busca-se uma função
$f(\mathbf{x})$ que generalize para transações não vistas. A saída é uma
probabilidade, convertida em decisão por um limiar $\tau$. A escolha de $\tau$
controla diretamente o compromisso entre Precisão e Recall.

\subsection{Aprendizado com classes desbalanceadas}

Quando uma classe é muito rara, algoritmos que minimizam o erro global tendem a
ignorá-la. Há duas famílias de soluções. Nas abordagens em nível de algoritmo,
atribui-se peso maior aos erros na classe minoritária (\emph{class weights}).
Nas abordagens em nível de dados, a distribuição de treino é reamostrada. O
SMOTE (\emph{Synthetic Minority Over-sampling Technique}) cria exemplos
sintéticos da classe minoritária por interpolação entre um exemplo e seus
vizinhos mais próximos \cite{chawla:2002}. A reamostragem deve ser aplicada
somente nos dados de treino, pois aplicá-la antes da divisão contamina o
conjunto de teste e superestima o desempenho. Além disso, ela altera as
probabilidades estimadas pelo modelo, o que pode exigir calibração posterior
\cite{dalpozzolo:2015}.

\subsection{Modelos de classificação}

A \textbf{Regressão Logística} modela a probabilidade de fraude como uma função
sigmoide de uma combinação linear dos atributos. É simples e interpretável,
por isso serve de referência. A \textbf{Random Forest} \cite{breiman:2001}
combina muitas árvores de decisão treinadas sobre amostras \emph{bootstrap} e
subconjuntos aleatórios de atributos. A votação reduz a variância e captura
relações não lineares. O \textbf{XGBoost} \cite{chen:2016} constrói árvores
sequencialmente por \emph{gradient boosting}, cada uma corrigindo os resíduos
das anteriores, com regularização e parâmetro próprio para desbalanceamento. As
\textbf{Máquinas de Vetores de Suporte} (SVM) \cite{cortes:1995} buscam o
hiperplano de margem máxima entre as classes, podendo usar \emph{kernels} para
fronteiras não lineares. Seu custo de treinamento cresce de forma
aproximadamente quadrática com o número de exemplos, o que exige subamostragem
ou versão linear em bases grandes.

\subsection{Métricas de avaliação}

Sejam VP, FP, VN e FN os verdadeiros positivos, falsos positivos, verdadeiros
negativos e falsos negativos. A \emph{Precisão} ($\frac{VP}{VP+FP}$) mede a
fração de alertas que são fraudes reais. O \emph{Recall}
($\frac{VP}{VP+FN}$) mede a fração de fraudes detectadas, e o \emph{F1} é a
média harmônica das duas. A área sob a curva Precisão--Recall (PR-AUC, ou
\emph{average precision}) resume o desempenho em todos os limiares e é mais
informativa que a área sob a curva ROC quando a classe positiva é rara
\cite{saito:2015}. O coeficiente de correlação de Matthews (MCC) considera as
quatro células da matriz de confusão e só é alto quando o classificador acerta
bem ambas as classes \cite{chicco:2020}.

\subsection{Trabalhos relacionados}

Em \cite{bhattacharyya:2011}, os autores compararam Regressão
Logística, SVM e Random Forest em dados reais de transações de cartão e
concluíram que a Random Forest apresentou o melhor desempenho geral, destacando
a importância de avaliar os modelos sob diferentes proporções de fraude.
Os autores do dataset utilizado neste trabalho
\cite{dalpozzolo:2015} mostraram que a subamostragem da classe majoritária
distorce as probabilidades previstas e propuseram uma correção de calibração.
Em trabalho posterior, \cite{dalpozzolo:2018}
argumentam que avaliações realistas devem considerar a ordem temporal das
transações, o atraso na chegada dos rótulos e o número limitado de alertas que
analistas conseguem verificar, propondo métricas de precisão sobre os alertas.

No mesmo dataset, \cite{awoyemi:2017} avaliaram Naïve
Bayes, k-vizinhos mais próximos (KNN) e Regressão Logística com amostragem
híbrida, e o KNN obteve o melhor desempenho entre os três.
Já \cite{varmedja:2019} aplicaram SMOTE e compararam
Regressão Logística, Random Forest, Naïve Bayes e Perceptron Multicamadas, com
melhor resultado geral para a Random Forest. Ambos os estudos, porém,
aplicaram a reamostragem com pouco detalhe sobre o protocolo e enfatizaram a
acurácia, o que dificulta a comparação direta.
O estudo de \cite{carcillo:2021} mostrou que incluir escores de
detecção de anomalias não supervisionada como atributos de modelos
supervisionados pode melhorar a detecção. Por fim, o manual de
\cite{leborgne:2022} sistematiza boas práticas
reprodutíveis, com ênfase em validação temporal e métricas voltadas a alertas.

Esses trabalhos fundamentam as decisões deste projeto: comparar modelos de
\emph{ensemble} (que se destacam na literatura) contra uma referência linear;
aplicar SMOTE apenas no treino; priorizar PR-AUC e MCC em vez de acurácia; e
incluir uma validação temporal. A Tabela~\ref{tab:relacionados} resume a
comparação.

\begin{table}[ht]
\centering
\caption{Síntese dos trabalhos relacionados}
\label{tab:relacionados}
\footnotesize
\begin{tabularx}{\textwidth}{@{}>{\raggedright\arraybackslash}p{3.3cm}>{\raggedright\arraybackslash}p{3.2cm}X@{}}
\toprule
\textbf{Trabalho} & \textbf{Técnicas} & \textbf{Solução obtida / relação com este projeto} \\
\midrule
Bhattacharyya et al. (2011) & RL, SVM, RF & RF com melhor desempenho geral; motiva o uso de \emph{ensembles}. \\
Dal Pozzolo et al. (2015) & Subamostragem, calibração & Correção das probabilidades após reamostragem; alerta para efeitos do balanceamento. \\
Dal Pozzolo et al. (2018) & Aprendizado com \emph{feedback} e rótulos atrasados & Avaliação realista e temporal; motiva a validação temporal. \\
Awoyemi et al. (2017) & NB, KNN, RL + amostragem híbrida & KNN superior no mesmo dataset; ênfase em acurácia. \\
Varmedja et al. (2019) & RL, RF, NB, MLP + SMOTE & RF com melhor resultado geral no mesmo dataset. \\
Carcillo et al. (2021) & Não supervisionado + supervisionado & Escores de anomalia como atributos melhoram a detecção; possível extensão futura. \\
Le Borgne et al. (2022) & Diversos & Boas práticas reprodutíveis, validação temporal e métricas voltadas a alertas. \\
\bottomrule
\end{tabularx}
\end{table}

% =====================================================================
\section{Metodologia e Resultados Esperados}

\subsection{Abordagem e etapas}

O projeto segue um fluxo típico de ciência de dados, implementado em Python com
scikit-learn \cite{pedregosa:2011} e imbalanced-learn \cite{lemaitre:2017}:

\begin{enumerate}
\item \textbf{Aquisição e análise exploratória} (N1): estrutura, qualidade,
desbalanceamento e separabilidade dos atributos.
\item \textbf{Preparação} (N1): remoção de duplicatas, engenharia de atributos,
divisão treino/teste e padronização.
\item \textbf{Modelos de referência} (N1): classificador trivial, Regressão
Logística (com pesos e com SMOTE) e Random Forest com pesos balanceados.
\item \textbf{Modelagem e ajuste} (N2): Random Forest, XGBoost e SVM, com busca
de hiperparâmetros por validação cruzada estratificada e comparação entre
pesos de classe e SMOTE.
\item \textbf{Escolha do limiar e avaliação} (N2): definição de $\tau$ em
validação, avaliação final no teste e na divisão temporal.
\item \textbf{Explicabilidade} (N2): importância por permutação e valores SHAP
\cite{lundberg:2017}.
\item \textbf{Protótipo} (N2): aplicação que recebe uma transação e devolve a
probabilidade de fraude e a decisão.
\end{enumerate}

Em todas as etapas, as transformações que aprendem parâmetros dos dados
(padronização e SMOTE) são encapsuladas em \emph{pipelines} ajustados somente
no treino (e em cada dobra da validação cruzada), evitando vazamento de
informação para o teste.

\subsection{Resultados esperados}

Espera-se que os modelos de \emph{ensemble} superem a Regressão Logística em
PR-AUC e MCC, em linha com a literatura \cite{bhattacharyya:2011,varmedja:2019},
e que SMOTE e pesos de classe aumentem o Recall à custa de parte da Precisão.
Espera-se também que o desempenho na divisão temporal seja igual ou inferior ao
da divisão aleatória, evidenciando o efeito da ordem temporal.

\subsection{Dataset}\label{sec:dataset}

\subsubsection{Origem e conteúdo}

Foi utilizado o dataset \emph{Credit Card Fraud Detection}, coletado pelo
Machine Learning Group da Université Libre de Bruxelles em colaboração com a
Worldline e disponibilizado publicamente no Kaggle
\cite{kaggle:ulb,dalpozzolo:2015}. Ele contém 284.807 transações realizadas em
dois dias de setembro de 2013 por portadores de cartão europeus, das quais 492
são fraudes (0,172\%). Os atributos \texttt{V1}--\texttt{V28} são componentes
principais (PCA) de atributos originais não divulgados por confidencialidade.
Apenas \texttt{Time} (segundos desde a primeira transação) e \texttt{Amount}
(valor) não foram transformados. A variável-alvo é \texttt{Class}.

\subsubsection{Qualidade, adequação e limitações}

Por serem transações reais rotuladas, os dados são representativos do
problema, e a anonimização por PCA elimina dados pessoais. Isso, porém, tem
custos: (i) \texttt{V1}--\texttt{V28} não têm significado semântico, o que
limita a engenharia de atributos e a explicabilidade das decisões para
usuários finais; (ii) como o horário real de início da coleta não é informado,
a hora derivada de \texttt{Time} é relativa; (iii) dois dias não permitem
avaliar sazonalidade nem \emph{concept drift}; (iv) as fraudes são poucas em
números absolutos, o que torna as métricas sensíveis a poucos erros e exige
cautela na interpretação; e (v) não há atributos demográficos, o que impede
auditar vieses contra grupos específicos. Essas limitações são consideradas na
interpretação dos resultados.

\subsubsection{Análise exploratória}

As 284.807 transações abrangem 48 horas. O conjunto não possui valores
ausentes, mas contém 1.081 linhas duplicadas, das quais 19 são fraudes. A
Figura~\ref{fig:classes} mostra o desbalanceamento: após a remoção das
duplicatas, restam 283.253 transações legítimas e 473 fraudes (0,167\%), uma
razão de aproximadamente 599:1.

\figura{fig01_distribuicao_classes.png}{0.5}{Distribuição das classes (escala
logarítmica)}{fig:classes}

Quanto ao valor (Figura~\ref{fig:amount}), as fraudes têm mediana
\emph{menor} que as transações legítimas (9,82 contra 22,00), mas média
\emph{maior} (123,87 contra 88,41). Isso indica uma concentração de fraudes em
valores muito baixos, possivelmente testes de validade do cartão, ao lado de
algumas fraudes de valor elevado. Transações de valor zero são 5,3\% das
fraudes e apenas 0,6\% das legítimas. O maior valor legítimo (25.691,16) supera
em muito a maior fraude (2.125,87), e a distribuição é fortemente assimétrica:
11,2\% das transações estão acima de $Q_3 + 1{,}5\,\mathrm{IQR}$. Por isso, optou-se pela
transformação logarítmica em vez da remoção de valores extremos, que em fraude
podem ser justamente os casos de interesse. A estatística de
Kolmogorov--Smirnov (KS) entre as classes para \texttt{Amount} foi de
0,265, indicando separação apenas moderada. O valor, isoladamente, não
identifica fraudes.

\figura{fig02_amount_por_classe.png}{0.6}{Distribuição de $\log(1+\mathrm{Amount})$
por classe}{fig:amount}

A Figura~\ref{fig:tempo} mostra o volume de transações e a taxa de fraude por
hora relativa. Há um ciclo nítido: entre as horas relativas 1 e 7 o volume
cai drasticamente (mínimo na hora 4), padrão compatível com o período noturno.
Nesse intervalo, a taxa de fraude sobe e atinge 1,45\% na hora 2, cerca de
5,5 vezes a média horária de 0,26\%. Como o denominador é pequeno nessas horas,
parte da variação pode ser ruído. Ainda assim, o padrão justifica incluir a
hora como atributo.

\figura{fig03_tempo_hora.png}{0.75}{Volume de transações e taxa de fraude por
hora relativa ao início da coleta}{fig:tempo}

Entre os componentes PCA, a correlação máxima em módulo entre pares foi de
0,019, confirmando a quase ortogonalidade esperada da PCA. Os atributos
mais correlacionados com a classe têm correlação negativa, com destaque para
\texttt{V17} ($-0{,}31$), \texttt{V14} ($-0{,}29$) e \texttt{V12} ($-0{,}25$)
(Figura~\ref{fig:corr}). Pela estatística KS, os mais separadores são
\texttt{V14} (0,84), \texttt{V10} (0,80), \texttt{V12} (0,78) e \texttt{V4}
(0,76) (Figura~\ref{fig:ks}), enquanto 4 componentes apresentaram KS inferior
a 0,1. Isso indica que a informação discriminativa
está concentrada em poucos componentes.

\figura{fig04_correlacao_classe.png}{0.6}{Atributos mais correlacionados com a
classe}{fig:corr}

\figura{fig05_top4_ks.png}{0.95}{Distribuições por classe dos quatro atributos
com maior estatística KS}{fig:ks}

\subsubsection{Preparação dos dados}

Com base na análise, foram tomadas as seguintes decisões: (i) remoção das
duplicatas antes da divisão, evitando que a mesma transação apareça no treino e
no teste; (ii) substituição de \texttt{Time} pela hora relativa em codificação
cíclica ($\sin$ e $\cos$ de $2\pi h/24$), que preserva a proximidade entre 23h
e 0h; (iii) transformação $\log(1+\mathrm{Amount})$; (iv) divisão estratificada
80/20 com semente fixa, resultando em 226.980 transações de treino
(378 fraudes) e 56.746 de teste (95 fraudes); (v) padronização
com \texttt{RobustScaler} (mediana e intervalo interquartil), menos sensível a
valores extremos; e (vi) SMOTE aplicado somente ao treino, dentro do
\emph{pipeline}. O conjunto final possui 31 atributos.

% =====================================================================
\section{Resultados Parciais}

\subsection{Aspectos éticos e responsabilidade no uso da IA}

\textbf{Privacidade e proteção de dados.} O dataset foi anonimizado pelos
autores antes da publicação e não contém dados pessoais identificáveis. O grupo
não fará nenhuma tentativa de reidentificação. Em uma aplicação real, dados de
transações são dados pessoais sujeitos à LGPD \cite{lgpd:2018}. O uso para
prevenção de fraudes precisaria de base legal adequada, minimização dos dados
coletados e controles de acesso e retenção.

\textbf{Custo assimétrico dos erros.} Um falso positivo bloqueia a compra
legítima de um cliente, o que pode ter consequências sérias (por exemplo, numa
compra de medicamento ou numa emergência). Um falso negativo transfere o
prejuízo ao cliente ou à instituição. Por isso, o projeto reporta os dois tipos
de erro separadamente, escolhe o limiar de decisão de forma explícita e propõe
que o modelo atue como \emph{apoio à decisão}: alertas vão para verificação (por
exemplo, confirmação com o cliente), e não para bloqueio definitivo automático.

\textbf{Transparência e explicabilidade.} A LGPD garante ao titular o direito de
solicitar a revisão de decisões tomadas unicamente com base em tratamento
automatizado (art.~20). Como \texttt{V1}--\texttt{V28} não têm significado
semântico, explicações em termos desses atributos são úteis para auditoria
técnica, mas não para o cliente. O projeto usará SHAP \cite{lundberg:2017}
para auditoria e reconhece essa limitação. Em produção, seriam necessários
atributos interpretáveis.

\textbf{Viés e equidade.} A ausência de atributos demográficos protege a
privacidade, mas impede verificar se o modelo erra mais para determinados
grupos, por exemplo, gerando mais bloqueios para perfis de consumo menos comuns.
Esse risco não pode ser descartado. Em um uso real, seria necessária auditoria
periódica das taxas de falso positivo por segmento.

\textbf{Robustez e responsabilidade.} Fraudadores adaptam seu comportamento, e
um modelo treinado em dois dias de 2013 não deve ser tratado como válido para o
cenário atual. Resultados obtidos neste dataset público não devem ser
apresentados como garantia de desempenho em produção. Medidas de mitigação
seriam o monitoramento contínuo, o retreinamento periódico e a supervisão
humana dos alertas.

\subsection{Resultados parciais obtidos}

A Tabela~\ref{tab:baselines} apresenta os resultados dos modelos de referência
no conjunto de teste, com limiar $\tau = 0{,}5$, e a PR-AUC média na validação
cruzada estratificada com 5 dobras no treino.

\begin{table}[ht]
\centering
\caption{Resultados dos modelos de referência no conjunto de teste}
\label{tab:baselines}
\footnotesize
\begin{tabular}{@{}lccccccc@{}}
\toprule
\textbf{Modelo} & \textbf{Acur.} & \textbf{Prec.} & \textbf{Recall} & \textbf{F1} & \textbf{MCC} & \textbf{PR-AUC} & \textbf{PR-AUC (CV)} \\
\midrule
Dummy (majoritária) & 0,998 & 0,000 & 0,000 & 0,000 & 0,000 & 0,002 & 0,002 $\pm$ 0,000 \\
RL (pesos) & 0,975 & 0,056 & \textbf{0,874} & 0,105 & 0,217 & 0,680 & 0,756 $\pm$ 0,027 \\
RL + SMOTE & 0,974 & 0,053 & \textbf{0,874} & 0,100 & 0,212 & 0,678 & 0,756 $\pm$ 0,025 \\
RF (pesos) & 0,9995 & \textbf{0,972} & 0,726 & \textbf{0,831} & \textbf{0,840} & \textbf{0,801} & \textbf{0,837} $\pm$ 0,030 \\
\bottomrule
\end{tabular}
\end{table}

O classificador trivial confirma que a acurácia é enganosa neste problema:
ele atinge 99,8\% sem detectar nenhuma fraude. As métricas focadas na classe
positiva revelam diferenças que a acurácia esconde.

As duas variantes de Regressão Logística detectaram 83 das 95 fraudes do teste
(Recall de 87,4\%), mas geraram cerca de 1.400 falsos positivos
(Figura~\ref{fig:cm}), ou seja, aproximadamente 17 alertas falsos para cada
fraude detectada (Precisão de 5,6\%). A Random Forest apresentou o
comportamento oposto: apenas 2 falsos positivos (Precisão de 97,2\%), mas 26
fraudes não detectadas (Recall de 72,6\%). Nas métricas que resumem o
equilíbrio entre os erros, a Random Forest foi claramente superior: MCC de 0,84
contra 0,22 e PR-AUC de 0,80 contra 0,68. O resultado da validação cruzada
(0,84 $\pm$ 0,03 contra 0,76 $\pm$ 0,03) confirma que a diferença não se deve à
partição de teste escolhida.

Dois achados merecem destaque. Primeiro, o SMOTE não trouxe ganho em relação
aos pesos de classe na Regressão Logística: os resultados foram praticamente
idênticos, com ligeiro aumento de falsos positivos (1.477 contra 1.406) e maior
custo computacional. Na N2, a comparação entre as duas estratégias será
repetida nos modelos de árvores. Segundo, a Random Forest teve ROC-AUC
\emph{menor} que a Regressão Logística (0,935 contra 0,963), embora tenha PR-AUC
muito maior (Figura~\ref{fig:pr}). Isso ilustra empiricamente o argumento de
\cite{saito:2015}: com classes raras, a curva ROC é dominada pela grande massa
de verdadeiros negativos e pode ordenar os modelos de forma diferente da curva
Precisão--Recall, que reflete melhor o custo operacional dos alertas.

Os perfis opostos de erro mostram também que o limiar padrão $\tau = 0{,}5$ não
é adequado para ambos. A escolha do limiar em validação, conforme o custo
relativo de falsos positivos e falsos negativos, é portanto uma etapa
necessária da N2.

\figura{fig06_curvas_pr.png}{0.6}{Curvas Precisão--Recall no conjunto de
teste}{fig:pr}

\figura{fig07_matrizes_confusao.png}{0.95}{Matrizes de confusão dos modelos de
referência}{fig:cm}

A importância dos atributos na Random Forest (Figura~\ref{fig:imp}) é
coerente com a análise exploratória. Os seis atributos mais importantes
(\texttt{V14}, \texttt{V10}, \texttt{V12}, \texttt{V4}, \texttt{V11} e
\texttt{V17}) são exatamente os seis de maior estatística KS, e \texttt{V14}
concentra sozinho 22\% da importância total. Os atributos derivados de valor e
hora não aparecem entre os dez mais importantes, o que sugere que sua
informação é, em grande parte, redundante com os componentes PCA.

\figura{fig08_importancia_rf.png}{0.6}{Atributos mais importantes na Random
Forest}{fig:imp}

Na divisão temporal (treino nas primeiras 80\% das transações e teste nas 20\%
finais, com 74 fraudes), a Random Forest manteve PR-AUC de 0,80 e Precisão de
98,0\%, mas o Recall caiu de 72,6\% para 67,6\%. A Regressão Logística teve
PR-AUC de 0,76, mas dobrou os falsos positivos (2.836) com o limiar padrão. Ao
contrário do esperado, não houve degradação expressiva da Random Forest. Como
o conjunto cobre apenas dois dias, isso não demonstra robustez a mudanças de
padrão de fraude ao longo de meses: indica apenas que não há mudança forte
dentro dessa janela.

Uma limitação deve ser considerada ao interpretar todos esses números: com
95 fraudes no teste, cada fraude corresponde a cerca de 1,05 ponto percentual
de Recall. Por isso, as comparações devem se apoiar também na validação
cruzada. Esses resultados estabelecem o patamar que os modelos da próxima etapa
(Random Forest ajustada, XGBoost e SVM) devem superar, a saber PR-AUC de 0,80
e MCC de 0,84, com atenção especial ao aumento do Recall sem perda
significativa de Precisão.

\subsection{Repositório}

Os artefatos deste trabalho (artigo, descrição do dataset, notebook e
resultados) estão disponíveis em:
\pend{\texttt{https://github.com/USUARIO/REPOSITORIO}}.

\bibliographystyle{sbc}
\bibliography{referencias}

\end{document}
